Because the interaction of the nuclear spins with the electrons is exceedingly weak, the probability of changing $I$ is very small even in molecular collisions. Molecules that differ in the parity of $I$, and accordingly possess only symmetric or only antisymmetric terms, therefore behave in practice as distinct modifications of a substance. Examples are the so-called orthohydrogen and parahydrogen: in a molecule of the former, the spins $i=1/2$ of the two nuclei are parallel ($I=1$), whereas in one of the latter they are antiparallel ($I=0$).
